\chapter{Short-Term Interest-Rate Futures}\label{ch:m2:short-term-interest-rate-futures}

A trader looks at the December 2026 one-month SOFR future. It trades at 96.040,
which says that the overnight rate will average 3.960\% over the month. The
Federal Reserve meets on 9 December; the November contract says the rate will
be 3.8975\% going into the meeting. Nine days of December at 3.8975\% and
twenty-two at an unknown rate must average 3.960\%: the unknown rate is
3.986\%, 8.8 basis points higher, and if the only choices are to hold or to
raise by a quarter point, the market is pricing a 35\% chance of a hike. The
trader's desk thinks the chance is lower, but before selling the contract it
has to ask one more question, about the last day of December, when the
overnight rate is usually higher for reasons that have nothing to do with the
Federal Reserve. The prices in this chapter are illustrative; the meeting dates
and the contracts' rules are real. It shows how futures on overnight rates
turn a central bank's calendar into prices, and how to read them back.

\section{Overnight-rate futures}

\begin{definition}[Overnight-rate future]\label{def:m2:short-term-interest-rate-futures:future}
An \emph{overnight-rate future}\index{overnight-rate future} is a cash-settled
futures contract whose final price is 100 minus a rate computed from the daily
fixings of an overnight benchmark (\cref{def:m2:central-banks-and-the-short-rate:rfr})
over a reference period: their arithmetic average over a calendar month for
one-month contracts, their compounded rate over a quarter for three-month
contracts. Before expiry the price is the market's expectation of that
settlement, adjusted for convexity (\cref{sec:m2:short-term-interest-rate-futures:convexity}).
\end{definition}

Quoting $100 - R$ makes the future behave like a bond: it rises when rates
fall. Each contract has a fixed value per basis point, so a position's
interest-rate risk is simply its number of contracts times that value, and
gains and losses are paid daily as variation margin (One Quant Book 1,
chapter 5).

\begin{dated}{2026-09}{Contracts on overnight rates}\label{dat:m2:short-term-interest-rate-futures:contracts}
\textbf{CME three-month SOFR} (\texttt{SR3}): USD~2\,500 times the index, USD~25 per
basis point; 100 minus SOFR compounded over the \omterm{def:m2:short-term-interest-rate-futures:imm}{reference quarter}, which runs
from the third Wednesday of the third month before the delivery month to the
third Wednesday of the delivery month; tick a quarter of a basis point for the
nearest contracts, half otherwise; 39 quarterly contracts, nearly ten years.
\textbf{CME one-month SOFR} (\texttt{SR1}): USD~4\,167 times the index, USD~41.67 per
basis point; 100 minus the arithmetic average of SOFR over the calendar month.
\textbf{ICE three-month SONIA}: GBP~2\,500 times the index; compounded SONIA from
the third Wednesday of the delivery month to the next quarterly one; 25
quarterly contracts. \textbf{Euribor}, the euro's surviving interbank rate,
is published at 11:00 CET for five tenors by its administrator, and futures
on three-month Euribor trade alongside those on overnight rates.
\end{dated}

\section{IMM dates, packs, bundles and strips}

\begin{definition}[IMM date and reference quarter]\label{def:m2:short-term-interest-rate-futures:imm}
An \emph{IMM date}\index{IMM date} is the third Wednesday of March, June,
September or December. A three-month contract's \emph{reference
quarter}\index{reference quarter} is the period between two consecutive IMM
dates over which its rate is compounded.
\end{definition}

The dates are used beyond futures: forward-rate agreements and swaps are
often dated to them, so that their hedges line up with the \omterm{def:m2:short-term-interest-rate-futures:strip}{futures strip}.

\begin{definition}[Futures strip, pack, bundle]\label{def:m2:short-term-interest-rate-futures:strip}
A \emph{futures strip}\index{futures strip} is a sequence of consecutive
quarterly contracts, whose rates together describe the path of the short rate.
A \emph{pack}\index{pack} is four consecutive quarterly contracts traded in
equal amounts as one instrument, a year of the strip; a
\emph{bundle}\index{bundle} is the first $n$ years of consecutive contracts
traded together, from the front.
\end{definition}

\begin{proposition}[A strip is a term rate]\label{prop:m2:short-term-interest-rate-futures:strip}
If consecutive reference periods of lengths $\delta_1, \dots, \delta_n$ (in
years) have futures rates $R_1, \dots, R_n$ (convexity ignored), a deposit
spanning them all has the rate
\[
R_{1\dots n} \;=\; \frac{\prod_{i}(1 + R_i\,\delta_i) - 1}{\sum_i \delta_i},
\]
and buying the strip locks in that rate for an investor who will roll a
deposit through the periods.
\end{proposition}
\begin{proof}
One unit invested at each period's rate grows to the product; the strip's
daily settlements replicate the difference between the realised and the
locked rates, period by period, to first order.
\end{proof}

\begin{omfigure}
\begin{tikzpicture}[font=\footnotesize,x=1.05cm]
\draw[thick,-{Stealth[length=2.5mm]}] (0,0) -- (12.9,0) node[right]{time};
\foreach \x/\l in {0/1 Sep,1/1 Oct,2/1 Nov,3/1 Dec,4/1 Jan,5/1 Feb,6/1 Mar}{\draw (\x*1.9,0.1) -- (\x*1.9,-0.1) node[below]{\l};}
% one-month contracts
\foreach \x/\l in {1/Oct,2/Nov,3/Dec,4/Jan}{\draw[omDef,line width=2.5pt] (\x*1.9+0.05,0.45) -- (\x*1.9+1.85,0.45); \node[above,text=omDef] at (\x*1.9+0.95,0.5) {SR1 \l};}
% x = 1.9 x (months after 1 Sep + day fraction): 16 Sep 0.50, 16 Dec 3.48, 17 Mar 6.52
\draw[omThm,line width=2.5pt] (0.50*1.9,1.35) -- (3.48*1.9,1.35);
\node[above,text=omThm] at (2.0*1.9,1.4) {SR3 December: 16 Sep -- 16 Dec};
\draw[omProp,line width=2.5pt] (3.48*1.9,1.35) -- (6.52*1.9,1.35);
\node[above,text=omProp] at (5.0*1.9,1.4) {SR3 March: 16 Dec -- 17 Mar};
% FOMC decisions
\foreach \x/\l in {0.50/16 Sep,1.87/28 Oct,3.26/9 Dec,4.84/27 Jan}{\draw[omThm,thick,dashed] (\x*1.9,-0.6) -- (\x*1.9,0.3); \node[below,text=omThm] at (\x*1.9,-0.6) {\l};}
\node[anchor=east,text=omThm] at (-0.1,-0.8) {FOMC:};
\end{tikzpicture}

\omcaption{Reference periods of the contracts of this chapter, September 2026
to March 2027. One-month contracts average their calendar month; three-month
contracts compound from \omterm{def:m2:short-term-interest-rate-futures:imm}{IMM date} to \omterm{def:m2:short-term-interest-rate-futures:imm}{IMM date}, and are named after the month
their period \emph{ends}. The dashed lines are the Federal Open Market
Committee's decision days; each new rate applies from the next
day.}\label{fig:m2:short-term-interest-rate-futures:calendar}
\end{omfigure}

\section{Pricing a meeting}

\begin{definition}[Implied policy path]\label{def:m2:short-term-interest-rate-futures:path}
The \emph{implied policy path}\index{implied policy path} is the sequence of policy
rates the market expects after each scheduled meeting of the central bank, read
from the prices of short-term interest-rate futures (or overnight index swaps)
as if those prices were expectations, any risk premium neglected.
\end{definition}

\begin{method}[Reading a meeting from a one-month contract]\label{met:m2:short-term-interest-rate-futures:meeting}
Let a month of $N$ calendar days have a decision on day $m$, the new rate
applying from day $m + 1$. Let $\bar R$ be the rate implied by the month's
contract, $100 - F$, and $r_0$ the rate expected before the meeting (from the
previous month's contract or the current fixing). Remove any known \omterm{def:m2:money-markets:turn}{turn
premium} $\tau$ on $k$ days, then
\[
r_1 \;=\; \frac{N\bar R - k\tau - m\,r_0}{N - m},
\qquad
p \;=\; \frac{r_1 - r_0}{\Delta},
\]
where $p$ is the implied probability of a move of size $\Delta$ if the only
outcomes are no move and $\Delta$. Chain the months to obtain the rate expected
after each meeting.
\end{method}

\begin{example}[Three meetings]\label{ex:m2:short-term-interest-rate-futures:path}
From a rate of 3.87\% after the September 2026 decision, with illustrative
prices of 96.1275 (October), 96.1025 (November), 96.040 (December) and 95.990
(January 2027): the November contract implies 3.8975\% after the October
meeting, an 11\% chance of a hike; December implies 3.9856\% after the 9
December meeting, a 35\% chance; January implies 4.175\% after the 27 January
meeting (\cref{fig:m2:short-term-interest-rate-futures:path}).
\end{example}

\begin{remark}[Late meetings are fragile]\label{rem:m2:short-term-interest-rate-futures:fragile}
The multiplier $N/(N - m)$ turns a price error into an error in $r_1$. The
January meeting falls on the 27th, so only four days of the month carry the
new rate: an error of one basis point in the January price moves the implied
post-meeting rate by $31/4 = 7.75$ basis points, a third of a hike. The October
meeting is worse (three days). For late meetings desks use the next month's
contract, in which the new rate applies all month, as the example did for
October.
\end{remark}

\begin{omfigure}
\begin{tikzpicture}
\begin{axis}[width=11cm,height=4.8cm,
  ylabel={implied overnight rate (\%)},
  tick label style={font=\small},label style={font=\small},
  xmin=10,xmax=153,ymin=3.82,ymax=4.22,grid=major,grid style={gray!25},
  xtick={30,61,91,122,153},xticklabels={1 Oct,1 Nov,1 Dec,1 Jan,1 Feb},
  table/col sep=comma]
\addplot[omDef,very thick,const plot] table[x=t,y=rate]{figdata/markets-2/08-short-term-interest-rate-futures/path.csv};
\draw[omThm,thick,dashed] (axis cs:57,3.82) -- (axis cs:57,4.22);
\draw[omThm,thick,dashed] (axis cs:99,3.82) -- (axis cs:99,4.22);
\draw[omThm,thick,dashed] (axis cs:148,3.82) -- (axis cs:148,4.22);
\node[font=\footnotesize,anchor=north west,text=omThm] at (axis cs:101,3.975) {9 Dec: 35\%};
\node[font=\footnotesize,anchor=east,text=omThm] at (axis cs:146,4.12) {27 Jan};
\end{axis}
\end{tikzpicture}

\omcaption{The path of the overnight rate implied by the illustrative one-month
contracts of \cref{ex:m2:short-term-interest-rate-futures:path}, stepping at
each decision (dashed). The steps are expectations, averages over hold and
hike, not forecasts of a quarter-point move: 35\% of 25 basis points is the 8.8
of December. Data: the chapter's tutorial.}\label{fig:m2:short-term-interest-rate-futures:path}
\end{omfigure}

\begin{omfigure}
\begin{tikzpicture}
\begin{axis}[width=11cm,height=4.4cm,
  xlabel={assumed turn premium on 31 December (bp)},ylabel={hike probability (\%)},
  tick label style={font=\small},label style={font=\small},
  xmin=0,xmax=40,ymin=26,ymax=37,grid=major,grid style={gray!25},
  table/col sep=comma]
\addplot[omThm,very thick] table[x=turn,y=prob]{figdata/markets-2/08-short-term-interest-rate-futures/turn.csv};
\end{axis}
\end{tikzpicture}

\omcaption{The December hike probability implied by the same price as the
year-end turn assumed on the 31 December fixing grows: 35.2\% with no turn, 33.4\%
with 10 basis points, 0.18 percentage points less for each basis point of turn.
The turn is a money-market fact (\cref{ch:m2:money-markets}) that a policy
reading must remove first. Data: the chapter's weekend
problem.}\label{fig:m2:short-term-interest-rate-futures:turn}
\end{omfigure}

\section{The convexity adjustment}\label{sec:m2:short-term-interest-rate-futures:convexity}

A future is settled every day; a forward rate agreement or a swap is settled
once. The difference matters because the gains of a short futures position
(which profits when rates rise) arrive when rates are high and can be
reinvested at high rates, and its losses arrive when rates are low and can be
financed cheaply. The short enjoys this, so it accepts a futures rate higher
than the forward rate.

\begin{definition}[Convexity adjustment]\label{def:m2:short-term-interest-rate-futures:convexity}
The \emph{convexity adjustment}\index{convexity adjustment} of a short-rate
future is the difference between its implied rate and the forward rate for
the same period.
\end{definition}

\begin{proposition}[The adjustment in a Gaussian model]\label{prop:m2:short-term-interest-rate-futures:holee}
If the short rate follows $dr = \theta(t)\,dt + \sigma\,dW$ (the Ho--Lee model),
the futures rate for a period $[T_1, T_2]$ exceeds the forward rate by
approximately $\tfrac12\,\sigma^2\,T_1\,T_2$. One Quant Book 6 derives it.
\end{proposition}
\admitted

\begin{example}[Five and ten years out]\label{ex:m2:short-term-interest-rate-futures:convexity}
With a normal volatility of 1\% a year, the three-month contract starting in
five years carries $\frac12 \times 0.01^2 \times 5 \times 5.25 = 13.1$ basis points of
adjustment; the one starting in ten years, 51.25. At the front it is negligible;
at the back of the strip it is larger than a typical move in a day, and a
curve built from futures without it is wrong
(\cref{fig:m2:short-term-interest-rate-futures:convexity}).
\end{example}

\begin{omfigure}
\begin{tikzpicture}
\begin{axis}[width=11cm,height=4.6cm,
  xlabel={start of the three-month period (years)},ylabel={adjustment (bp)},
  tick label style={font=\small},label style={font=\small},
  xmin=0,xmax=10,ymin=0,ymax=78,grid=major,grid style={gray!25},
  legend columns=3,legend style={font=\footnotesize,at={(0.5,-0.3)},anchor=north,draw=none,fill=none,/tikz/every even column/.append style={column sep=8pt}},
  table/col sep=comma]
\addplot[omDef,thick,dashed] table[x=t1,y=s08]{figdata/markets-2/08-short-term-interest-rate-futures/convexity.csv};
\addplot[omThm,very thick] table[x=t1,y=s10]{figdata/markets-2/08-short-term-interest-rate-futures/convexity.csv};
\addplot[omProp,thick,dotted] table[x=t1,y=s12]{figdata/markets-2/08-short-term-interest-rate-futures/convexity.csv};
\legend{$\sigma = 0.8\%$,$\sigma = 1.0\%$,$\sigma = 1.2\%$}
\end{axis}
\end{tikzpicture}

\omcaption{The Ho--Lee \omterm{def:m2:short-term-interest-rate-futures:convexity}{convexity adjustment} of a three-month futures rate by
the start of its period, for three volatilities. It grows with the square of
the horizon and of the volatility: negligible for the first contracts, tens of
basis points for the last. Data: the chapter's tutorial.}\label{fig:m2:short-term-interest-rate-futures:convexity}
\end{omfigure}

\section{Tutorial: extracting a policy path}\label{tut:m2:short-term-interest-rate-futures:path}

\begin{tutorial}
\textbf{Goal.} Read the rate expected after each meeting from one-month futures,
with and without a year-end turn, and compute strip rates and \omterm{def:m2:short-term-interest-rate-futures:convexity}{convexity
adjustments}. \textbf{End state:}
\cref{fig:m2:short-term-interest-rate-futures:path,fig:m2:short-term-interest-rate-futures:turn}
and the numbers of \cref{ex:m2:short-term-interest-rate-futures:path}.
\begin{tutsteps}
\item \textbf{One month.} Remove the turn, then solve for the post-meeting rate.
\omcode{code/firm/meetings/firm_meetings.py}{32}{39}{The rate after a meeting from the month's average.}\label{lst:m2:short-term-interest-rate-futures:after}
\item \textbf{Chain the months}, and turn changes into probabilities.
\omcode{code/firm/meetings/firm_meetings.py}{42}{53}{The policy path, and the probability of a move.}\label{lst:m2:short-term-interest-rate-futures:path}
\item \textbf{Strips.}
\omcode{code/firm/meetings/firm_meetings.py}{56}{62}{The rate of a strip of consecutive periods.}\label{lst:m2:short-term-interest-rate-futures:strip}
\item \textbf{Run} \texttt{stir\_demo.path()} and \texttt{stir\_demo.december(10)}:
3.9856\% and 35.2\% without a turn, 3.9810\% and 33.4\% with 10 basis points.
\end{tutsteps}
\textbf{What to change next.} Replace the one-month contracts by three-month
contracts compounding over IMM quarters, with two meetings inside one quarter,
and solve for both (\cref{exo:m2:short-term-interest-rate-futures:7}).
\end{tutorial}

\section{Build: the policy-path extractor}\label{bld:m2:short-term-interest-rate-futures:meetings}

\begin{build}
\bfield{Purpose} The miniature firm's rates traders, and its curve builder
(\cref{ch:m2:interest-rate-swaps}), need the market's expected policy path:
the traders to take views meeting by meeting, the curve to put its steps on
meeting dates rather than smoothing through them.
\bfield{Interface} \texttt{MonthContract(year, month, price, meeting,
turn\_days, turn\_bp)}; \texttt{rate\_after(contract, rate\_before)};
\texttt{policy\_path(contracts, start\_rate)};
\texttt{move\_probability(after, before, step)};
\texttt{strip\_rate(rates, fractions)};
\texttt{convexity\_adjustment(sigma, t1, t2)}.
\bfield{Rules} Calendar-day averages; the new rate applies from the day after
a decision; a \omterm{def:m2:money-markets:turn}{turn premium} removed before solving; rates in percent.
\bfield{Acceptance tests} \texttt{code/firm/meetings/tests/}: a month without
a meeting returns its average; a constructed path is recovered exactly; a
turn is removed; chaining; strip compounding; the adjustment's value at five
years.
\bfield{Stretch} Three-month contracts (compounding, several meetings per
quarter); the probability of 25 against 50 basis points from two adjacent
contracts; the series of implied paths over time as a data product.
\end{build}

\begin{omsources}
\item CME Group, Rulebook Chapters 460 and 461 (three-month and one-month SOFR
futures); ICE, \emph{Three Month SONIA Index Futures}; EMMI, \emph{Euribor}.
\item Federal Reserve, \emph{FOMC meeting calendars}, 2026 and 2027.
\item T.\,S.\,Y.\ Ho and S.-B.\ Lee, ``Term structure movements and pricing
interest rate contingent claims'', \emph{Journal of Finance} 41, 1986.
\item G.\ Burghardt, \emph{The Eurodollar Futures and Options Handbook},
McGraw-Hill, 2003, for the market practice the SOFR contracts inherited.
\end{omsources}

\section{Exercises}

\begin{exercise}[$\star$]\label{exo:m2:short-term-interest-rate-futures:1}
A three-month SOFR contract trades at 95.965. What rate does it imply? It rises
3.5 basis points: how many half-basis-point ticks, and what is the gain on 200
contracts?
\end{exercise}

\begin{exercise}[$\star$]\label{exo:m2:short-term-interest-rate-futures:2}
Give the four \omterm{def:m2:short-term-interest-rate-futures:imm}{IMM dates} of 2027 and the \omterm{def:m2:short-term-interest-rate-futures:imm}{reference quarter} of the June 2027
three-month SOFR contract.
\end{exercise}

\begin{exercise}[$\star$]\label{exo:m2:short-term-interest-rate-futures:3}
The November one-month contract trades at 96.1025. What average rate does it
imply? What does a 2-basis-point rise in that rate cost a holder of 100 long
contracts?
\end{exercise}

\begin{exercise}[$\star\star$]\label{exo:m2:short-term-interest-rate-futures:4}
The October meeting is on the 28th. By how many basis points does a
1-basis-point error in the October price move the implied post-meeting rate?
Which contract would you use instead, and why?
\end{exercise}

\begin{exercise}[$\star\star$]\label{exo:m2:short-term-interest-rate-futures:5}
Four consecutive quarterly contracts imply 3.95\%, 4.05\%, 4.10\% and 4.12\% over
periods of 91 days each (actual/360). Give the one-year strip rate.
\end{exercise}

\begin{exercise}[$\star\star$]\label{exo:m2:short-term-interest-rate-futures:6}
A trader hedges a sterling position with the ``December'' three-month SONIA
future, believing it covers the same period as the ``December'' three-month
SOFR future. What is wrong?
\end{exercise}

\begin{exercise}[$\star\star\star$]\label{exo:m2:short-term-interest-rate-futures:7}
\emph{Coding.} With \texttt{convexity\_adjustment}, give the adjustment of the
three-month contract starting in five years and in ten years at 1\% volatility.
By how much does it change if volatility rises to 1.2\%?
\end{exercise}

\begin{exercise}[$\star\star\star$]\label{exo:m2:short-term-interest-rate-futures:8}
\emph{Find the flaw.} ``The October contract implies a 10\% chance of a hike on
28 October and the November contract 11\%: the market is internally
inconsistent, and the difference is an arbitrage.'' Correct it.
\end{exercise}

\section{Problem: Hike or Hold}

\begin{problem}[{Weekend problem --- the December 2026 meeting}]\label{pb:m2:short-term-interest-rate-futures:1}
The November one-month SOFR future trades at 96.1025 and the December one at
96.0400. The Committee decides on 9 December; the new rate applies from 10
December. The only possible outcomes are no change and a quarter-point hike.
Each contract is worth USD~41.67 per basis point.

\medskip\noindent\textbf{Part I --- The reading.}
\begin{enumerate}
\item What rate does the market expect going into the meeting?
\item How many December days are at the old rate, and how many at the new?
\item Give the post-meeting rate implied by the December price.
\item Give the implied probability of a hike.
\item How many basis points of December price does one percentage point of
probability represent?
\end{enumerate}

\medskip\noindent\textbf{Part II --- The turn.}
\begin{enumerate}[resume]
\item The 31 December fixing is expected 10 basis points above the rest.
Recompute the post-meeting rate and the probability.
\item And for 20 basis points?
\item By how much does each basis point of assumed turn reduce the probability?
\item Why does the turn fall in the December contract and not in January's?
\item How would you estimate the turn from market prices rather than assume it?
\end{enumerate}

\medskip\noindent\textbf{Part III --- The trade.}
\begin{enumerate}[resume]
\item A desk that thinks a hike is unlikely buys 1\,000 December contracts at
96.040. What is the final price and its P\&L if the Committee holds?
\item And if it hikes?
\item What is the expected P\&L at the market's 35\% probability?
\item At what probability does the desk's trade break even, and why is that the
market's number?
\item What else, besides the decision, moves the December settlement?
\end{enumerate}

\medskip\noindent\textbf{Part IV --- Judgement.}
\begin{enumerate}[resume]
\item Why is ``35\% chance of a hike'' an approximation even with no turn?
\item Why might the risk-neutral probability differ from the Committee members'
own projections?
\item How would a surprise hike in October change the December reading?
\item State the \emph{named result}: the implied December hike probability,
without and with a 10-basis-point year-end turn.
\item In one sentence: what does a one-month rate future price?
\end{enumerate}
\end{problem}

\section{Interview questions}

\begin{interviewq}[$\star$ \iqroles{trader, researcher}]\label{iq:m2:short-term-interest-rate-futures:1}
How do you read the probability of a rate hike from futures prices?
\end{interviewq}

\begin{interviewq}[$\star$ \iqroles{trader, bank}]\label{iq:m2:short-term-interest-rate-futures:2}
Why do short-rate futures quote $100$ minus the rate?
\end{interviewq}

\begin{interviewq}[$\star\star$ \iqroles{researcher, trader}]\label{iq:m2:short-term-interest-rate-futures:3}
What is the \omterm{def:m2:short-term-interest-rate-futures:convexity}{convexity adjustment} between futures and forward rates, and which
way does it go?
\end{interviewq}

\begin{interviewq}[$\star\star$ \iqroles{trader}]\label{iq:m2:short-term-interest-rate-futures:4}
What is a pack, and why would a trader trade the second year's pack rather
than the front contracts?
\end{interviewq}

\begin{interviewq}[$\star\star$ \iqroles{developer, researcher}]\label{iq:m2:short-term-interest-rate-futures:5}
You build a service that publishes meeting-by-meeting implied rates every
minute. What goes wrong at month-ends, around meetings late in a month, and at
the year-end?
\end{interviewq}

\begin{interviewq}[$\star\star\star$ \iqroles{researcher}]\label{iq:m2:short-term-interest-rate-futures:6}
Futures imply a 35\% chance of a hike. Is that the probability you should use
to decide a trade? What would you need to convert it?
\end{interviewq}
